\section{总结与延伸}

\subsection{讲者的核心总结}

\begin{figure}[H]
\centering
\includegraphics[width=0.95\textwidth]{slides-images/slide-084.jpg}
\caption{课程总结要点\protect\footnotemark}
\end{figure}
\footnotetext{Slides 第85页（最后一页）。}

Manning 教授在课程结尾总结了以下要点：

\begin{enumerate}
    \item \textbf{反向传播 = 链式法则的高效实现}：前向传播计算函数值，反向传播计算梯度
    \item \textbf{理解底层原理很重要}：虽然现代框架（PyTorch 等）自动化了一切，但理解反向传播的数学有助于调试和设计新架构
    \item \textbf{梯度消失/爆炸}：后续讲解 RNN 时会看到，多层梯度连乘可能导致梯度指数级增长或衰减
    \item \textbf{实践中的平衡}：框架预定义了常用操作的 forward/backward，用户像拼积木一样组合即可
\end{enumerate}

\subsection{全课知识图谱}

\begin{center}
\begin{tikzpicture}[
    node distance=0.9cm,
    every node/.style={draw, rounded corners, minimum width=4cm, minimum height=0.7cm, font=\small, align=center},
    arrow/.style={-{Stealth[length=3mm]}, thick}
]
\node (nn) {神经网络基础};
\node (act) [below=of nn] {激活函数与非线性};
\node (sgd) [below=of act] {SGD 与梯度计算};
\node (mc) [below=of sgd] {矩阵微积分};
\node (bp) [below=of mc] {反向传播算法};
\node (impl) [below=of bp] {自动微分与框架实现};

\draw[arrow] (nn) -- (act);
\draw[arrow] (act) -- (sgd);
\draw[arrow] (sgd) -- (mc);
\draw[arrow] (mc) -- (bp);
\draw[arrow] (bp) -- (impl);

\node[draw=none, right=0.8cm of nn, text width=5.5cm, font=\scriptsize, align=left] {层、权重、偏置、矩阵表示};
\node[draw=none, right=0.8cm of act, text width=5.5cm, font=\scriptsize, align=left] {Sigmoid, tanh, ReLU, Swish, GELU};
\node[draw=none, right=0.8cm of sgd, text width=5.5cm, font=\scriptsize, align=left] {$\theta^{new} = \theta^{old} - \alpha\nabla J(\theta)$};
\node[draw=none, right=0.8cm of mc, text width=5.5cm, font=\scriptsize, align=left] {Jacobian、链式法则、$\delta$、Hadamard 积};
\node[draw=none, right=0.8cm of bp, text width=5.5cm, font=\scriptsize, align=left] {计算图、前向/反向传播、梯度求和};
\node[draw=none, right=0.8cm of impl, text width=5.5cm, font=\scriptsize, align=left] {forward/backward API、梯度检验};
\end{tikzpicture}
\end{center}

\subsection{关键 Takeaways}

\begin{importantbox}{五条核心原则}
\begin{enumerate}
    \item \textbf{非线性是必须的}：没有激活函数，再多层也只是一个线性变换
    \item \textbf{矩阵微积分 = 单变量微积分 + 矩阵}：Jacobian、链式法则在形式上与标量版完全类似
    \item \textbf{$\delta$（上游梯度）是核心概念}：它代表从损失函数回传到当前位置的``误差信号''，是多个参数梯度共享的部分
    \item \textbf{反向传播的效率来自复用}：按逆拓扑序一次遍历，每个中间梯度只计算一次
    \item \textbf{形状约定简化工程实现}：将梯度整形为与参数相同的形状，使 SGD 更新变为简单的减法
\end{enumerate}
\end{importantbox}

\subsection{拓展阅读}

\begin{itemize}
    \item CS224N 官方讲义（Lecture Notes）与矩阵微积分教程：\url{https://web.stanford.edu/class/cs224n/}
    \item Stanford Math 51 在线教材（线性代数与多变量微积分）：\url{http://web.stanford.edu/class/math51/textbook.html}
    \item Karpathy, ``Yes you should understand backprop''：\url{https://karpathy.medium.com/yes-you-should-understand-backprop-e2f06eab496b}
    \item Goodfellow et al., \textit{Deep Learning} Chapter 6: Deep Feedforward Networks
    \item Justin Johnson, ``Backpropagation for a Linear Layer'': \url{https://web.eecs.umich.edu/~justincj/teaching/eecs442/notes/linear-backprop.html}
    \item PyTorch 官方教程 -- Autograd: \url{https://pytorch.org/tutorials/beginner/blitz/autograd_tutorial.html}
\end{itemize}

\end{document}
